\documentclass[12pt]{article}

% -------------------------------------------------
% PAQUETES
% -------------------------------------------------

\usepackage[utf8]{inputenc}
\usepackage[T1]{fontenc}
\usepackage[spanish]{babel}

\usepackage{amsmath}
\usepackage{amssymb}
\usepackage{booktabs}
\usepackage{array}
\usepackage{graphicx}
\usepackage{float}
\usepackage{geometry}
\usepackage{xcolor}
\usepackage{enumitem}
\usepackage{titlesec}
\usepackage{setspace}

% -------------------------------------------------
% CONFIGURACIÓN DE PÁGINA
% -------------------------------------------------

\geometry{
    left=2.5cm,
    right=2.5cm,
    top=2.5cm,
    bottom=2.5cm
}

\setstretch{1.2}

% -------------------------------------------------
% FORMATO DE TÍTULOS
% -------------------------------------------------

\titleformat{\section}
{\large\bfseries}
{\thesection.}{0.5em}{}

\titleformat{\subsection}
{\normalsize\bfseries}
{\thesubsection.}{0.5em}{}

% -------------------------------------------------
% DOCUMENTO
% -------------------------------------------------

\begin{document}

% -------------------------------------------------
% PORTADA
% -------------------------------------------------

\begin{titlepage}

\centering

\vspace*{2cm}

{\Large \textbf{MACROPOL: Economía Aplicada y Políticas Públicas}} \\[1cm]

{\LARGE \textbf{Stress Testing en Sistemas de Pensiones}} \\[0.5cm]

{\Large Taller de Escenarios Macro-Financieros} \\[2cm]

{\large Módulo 3: Construcción de Escenarios Macro-Financieros} \\[1cm]

\vfill

\textbf{Duración del ejercicio:} 45 minutos \\[0.5cm]

\textbf{Modalidad:} Trabajo grupal \\[0.5cm]

\textbf{Horizonte de simulación:} 24 meses

\vfill

{\large \today}

\end{titlepage}

% -------------------------------------------------
% CONTENIDO
% -------------------------------------------------

\section*{Contexto}

Este ejercicio corresponde al módulo de construcción de escenarios macro-financieros 
aplicados al stress testing de sistemas de pensiones de capitalización individual.

Cada grupo deberá analizar la lógica de construcción de los escenarios y comprender
cómo distintos modelos cuantitativos producen trayectorias macroeconómicas diferentes.

\vspace{0.5cm}

\section*{Objetivo}

Comprender:

\begin{itemize}
    \item cómo se construyen escenarios macro-financieros;
    \item cómo interactúan las variables en cada metodología;
    \item cómo interpretar escenarios baseline y adversos;
    \item cómo distinguir narrativa económica de resultados cuantitativos.
\end{itemize}

\vspace{0.5cm}

\section*{Instrucciones Generales}

Cada grupo deberá:

\begin{enumerate}
    \item Analizar el escenario asignado.
    
    \item Explicar:
    
    \begin{itemize}
        \item cómo fue construido;
        \item qué supuestos utiliza;
        \item qué variables dominan el escenario;
        \item qué dinámica económica representa.
    \end{itemize}
    
    \item Presentar resultados en un máximo de 10 minutos.
\end{enumerate}

\vspace{0.5cm}

\section*{Estructura Obligatoria de la Presentación}

\subsection*{1. Construcción del Escenario}

Responder:

\begin{itemize}
    \item ¿Qué metodología se utilizó?
    \item ¿El escenario es estadístico, histórico o hipotético?
    \item ¿Las variables interactúan entre sí?
\end{itemize}

\subsection*{2. Supuestos}

Identificar claramente:

\begin{itemize}
    \item supuestos del modelo;
    \item supuestos narrativos;
    \item horizonte temporal;
    \item naturaleza de los shocks.
\end{itemize}

\subsection*{3. Interpretación Económica}

Analizar:

\begin{itemize}
    \item crecimiento;
    \item inflación;
    \item tasa interbancaria;
    \item depreciación cambiaria;
    \item EMBI.
\end{itemize}

Responder:

\begin{itemize}
    \item ¿Qué variable domina el escenario?
    \item ¿Existe persistencia del shock?
    \item ¿El escenario parece coherente?
    \item ¿Qué tipo de episodio económico representa?
\end{itemize}

\vspace{0.5cm}

% -------------------------------------------------
% GRUPOS
% -------------------------------------------------

\section*{Grupos de Trabajo}

% -------------------------------------------------

\subsection*{Grupo 1 — Baseline ARIMA}

Analizar el escenario:

\begin{itemize}
    \item \texttt{baseline\_arima}
\end{itemize}

Preguntas guía:

\begin{enumerate}
    \item ¿Cómo funciona un modelo ARIMA?
    \item ¿Por qué este enfoque es univariado?
    \item ¿Qué ventajas tiene para proyección baseline?
    \item ¿Qué limitaciones presenta?
\end{enumerate}

% -------------------------------------------------

\subsection*{Grupo 2 — Baseline VAR}

Analizar el escenario:

\begin{itemize}
    \item \texttt{baseline\_var}
\end{itemize}

Preguntas guía:

\begin{enumerate}
    \item ¿Cómo funciona un modelo VAR?
    \item ¿Qué significa interdependencia entre variables?
    \item ¿Qué aporta el VAR frente al ARIMA?
    \item ¿Qué dinámica macroeconómica captura mejor?
\end{enumerate}

% -------------------------------------------------

\subsection*{Grupo 3 — Monte Carlo VAR}

Analizar:

\begin{itemize}
    \item \texttt{montecarlo\_p05};
    \item \texttt{montecarlo\_p50};
    \item \texttt{montecarlo\_p95}.
\end{itemize}

Preguntas guía:

\begin{enumerate}
    \item ¿Cómo se construye una simulación Monte Carlo?
    \item ¿Qué representan los percentiles?
    \item ¿Qué variable presenta mayor incertidumbre?
    \item ¿Cómo cambia la dispersión a través del tiempo?
\end{enumerate}

% -------------------------------------------------

\subsection*{Grupo 4 — Escenario Adverso Histórico}

Analizar:

\begin{itemize}
    \item \texttt{adverso\_historico}
\end{itemize}

Índice de estrés utilizado:

\[
0.25 \times inflacion +
0.20 \times interbancaria +
0.20 \times deprec\_fx +
0.20 \times embi -
0.15 \times crecimiento
\]

Preguntas guía:

\begin{enumerate}
    \item ¿Cómo se identificó la peor ventana histórica?
    \item ¿Qué variable domina el estrés?
    \item ¿Qué episodio económico parece representar?
    \item ¿Qué ventajas y limitaciones tiene este enfoque?
\end{enumerate}

% -------------------------------------------------

\subsection*{Grupo 5 — Escenario Adverso Hipotético}

Analizar:

\begin{itemize}
    \item \texttt{adverso\_hipotetico}
\end{itemize}

Preguntas guía:

\begin{enumerate}
    \item ¿Cómo se construyó el shock hipotético?
    \item ¿Qué narrativa económica representa?
    \item ¿Los shocks son persistentes o transitorios?
    \item ¿El escenario parece plausible?
\end{enumerate}

\vspace{0.5cm}

% -------------------------------------------------
% PRODUCTO FINAL
% -------------------------------------------------

\section*{Producto Esperado}

Cada grupo deberá entregar:

\begin{enumerate}
    \item Un resumen ejecutivo de máximo una página.
    
    \item Un gráfico con las principales trayectorias del escenario.
    
    \item Una explicación clara de:
    
    \begin{itemize}
        \item metodología;
        \item supuestos;
        \item interpretación económica.
    \end{itemize}
\end{enumerate}

\vspace{0.5cm}

% -------------------------------------------------
% CRITERIOS
% -------------------------------------------------

\section*{Criterios de Evaluación}

\begin{table}[H]
\centering

\begin{tabular}{lc}

\toprule

\textbf{Criterio} & \textbf{Peso} \\

\midrule

Comprensión metodológica & 35\% \\

Interpretación económica & 35\% \\

Claridad de exposición & 30\% \\

\bottomrule

\end{tabular}

\end{table}

\vspace{1cm}

% -------------------------------------------------
% MENSAJE FINAL
% -------------------------------------------------

\section*{Mensaje Final}

\begin{quote}
\itshape
“El objetivo del ejercicio no es producir pronósticos exactos, 
sino comprender cómo distintas metodologías construyen 
escenarios macro-financieros y cómo deben interpretarse.”
\end{quote}

\vfill

\begin{center}
\textbf{MACROPOL — Economía Aplicada y Políticas Públicas}
\end{center}

% -------------------------------------------------

\end{document}